\documentclass[11pt, letterpaper]{article}
\usepackage[margin=1in]{geometry}
\usepackage{enumitem}
\usepackage{parskip}
\usepackage{titlesec}
\usepackage{booktabs}
\usepackage{graphicx}
\usepackage{pgfgantt}
\usepackage[hidelinks]{hyperref}

% Section formatting
\titleformat{\section}{\normalfont\large\bfseries}{\thesection}{1em}{}[\titlerule]
\titleformat{\subsection}{\normalfont\normalsize\bfseries}{\thesubsection}{1em}{}
\titleformat{\subsubsection}{\normalfont\normalsize\itshape}{\thesubsubsection}{1em}{}
\titlespacing{\section}{0pt}{12pt}{6pt}
\titlespacing{\subsection}{0pt}{8pt}{4pt}
\titlespacing{\subsubsection}{0pt}{6pt}{2pt}

\begin{document}

% ─── Title Page ──────────────────────────────────────────────────────────────
\begin{titlepage}
    \centering
    \vspace*{2cm}
    {\large University of Pennsylvania}\\[0.5em]
    {\large CIS 5680: Game Design Practicum}\\[3cm]
    {\LARGE \textbf{Game Design Document}}\\[1em]
    {\Huge \textbf{Umbra}}\\[0.5em]
    {\large A Light and Shadow Puzzle Game}\\[3cm]
    {\normalsize Group Members: [Name 1], [Name 2], [Name 3]}\\[1em]
    {\normalsize Instructor: Dr.\ Stephen Lane}\\[1em]
    {\normalsize Spring 2026}
    \vfill
\end{titlepage}

% ─── Table of Contents ───────────────────────────────────────────────────────
\tableofcontents
\newpage

% ─── 1. Executive Summary ────────────────────────────────────────────────────
\section{Executive Summary}

\textbf{Umbra} is a top-down 2D spatial puzzle game built in Unreal Engine~5 for PC
(Windows). Shadows are rendered continuously and update in real time as objects in the
scene move: every light source and blocking object casts live shadow geometry that
physically affects the scene. Players manipulate portable lanterns and pivotable
spotlights---dragging and rotating them freely with the mouse within designer-specified
boundaries---to reshape the shadow pattern across the environment, and toggle between
morning and evening sun angles to flip the direction of every directional shadow
simultaneously. Blocking objects (crates, pillars) slide along fixed trajectories to
designer-specified stop positions when the player interacts with them. Shadow regions are
physically solid---a gap becomes crossable only when a shadow-bridge is cast across it,
and displacing that bridge destroys the path. Each level is a self-contained puzzle with
a correct configuration of light sources and blocking objects that opens a walkable path
from start to exit.

Umbra targets fans of clean, mechanic-first puzzle games such as \textit{The Witness},
\textit{Portal}, and \textit{Stephen's Sausage Roll}. There is no combat, no time
pressure, and no randomness---every solution is deterministic and discoverable through
spatial reasoning alone. Two foundational rules (shadows are physically solid; the sun
has exactly two states) combine to produce a surprisingly deep puzzle space without
requiring the player to learn a large system. The game is accessible to players aged
12~and~up and designed for pick-up-and-play sessions of 5--15 minutes per level.

Development is structured in three phases across a three-person team: an \emph{Alpha}
prototype delivering one fully playable level with the core shadow mechanic and sun
toggle (due Wed.~Mar.~18); a \emph{Beta} expanding to multiple puzzle levels with
polished art, audio, and a hint system (due Wed.~Mar.~25); and a \emph{Final} release
incorporating a tutorial level, level-select screen, and gameplay/trailer videos (due
Mon.~Mar.~30 and Wed.~Apr.~1). The full prototype is expected to be completed in
approximately five weeks of part-time development.

% ─── 2. Game Design ──────────────────────────────────────────────────────────
\section{Game Design}

    \subsection{High Concept}
    % TODO: Paste/adapt the approved high concept text from High_Concept_1.tex.

    \subsection{Design Goals}

        \subsubsection{Main Design Features}

            \paragraph{Player Goals and Objectives}
            The player is an explorer navigating a world where light and shadow are as
            solid as stone. The goal of each self-contained level is to discover the
            correct positions and orientations of light sources, the correct positions
            of blocking objects along their fixed trajectories, and the correct sun state
            (morning or evening) that together cast a shadow pattern opening a walkable
            path from start to exit. There is no time pressure and no combat; the sole
            challenge is spatial reasoning.

            \paragraph{Main Rules and Procedures}
            \begin{itemize}[leftmargin=1.5em, itemsep=2pt]
                \item Shadows are continuous and computed in real time by Unreal
                      Engine~5's lighting system. Any region covered by a shadow pattern
                      is physically solid and traversable; lit regions are open space.
                      Bridges, platforms, and doorways built from shadow geometry appear
                      or vanish as the shadow pattern changes.
                \item \textbf{Light sources} (lanterns, spotlights) are freely moved or
                      rotated by clicking and dragging with the mouse. Each light source
                      has a designer-specified boundary that limits how far it can travel
                      or rotate. Shadow geometry updates continuously as the player drags.
                \item \textbf{Blocking objects} (crates, pillars) slide along fixed,
                      pre-defined trajectories and snap to designer-specified stop
                      positions when the player clicks and drags them. They cannot be
                      freely placed; each object has exactly the set of positions its
                      trajectory permits.
                \item Players may toggle the sun state between morning (sun from the
                      east) and evening (sun from the west) at any time; this
                      immediately recomputes every directional shadow cast by blocking
                      objects and the environment.
                \item The level is solved when the player's character reaches the exit
                      through a walkable path that exists in the current shadow
                      configuration.
            \end{itemize}

            \paragraph{Player Resources}
            \begin{itemize}[leftmargin=1.5em, itemsep=2pt]
                \item \textbf{Portable lanterns.} Freely dragged to any position within
                      their movement boundary by clicking and holding with the mouse.
                      Cast a circular point-light shadow pattern that updates in real
                      time as the lantern is repositioned.
                \item \textbf{Wall-mounted spotlights.} Freely rotated about their fixed
                      mount point by clicking and dragging. Emit a directed cone of
                      light; rotating the spotlight continuously reshapes the shadow
                      it casts behind any blocking object in its beam.
                \item \textbf{Blocking objects} (crates, pillars). Slid to any of their
                      trajectory's designated stop positions by clicking and dragging
                      along their fixed path. Reposition them to add, remove, or redirect
                      a shadow cast by the sun or a nearby light source.
                \item \textbf{Sun-angle toggle.} A single key/button input switches
                      between morning and evening, instantly remapping every directional
                      shadow in the scene and producing an entirely new traversable
                      shadow pattern from the same geometry.
            \end{itemize}

            \paragraph{Conflicts}
            \begin{itemize}[leftmargin=1.5em, itemsep=2pt]
                \item \textbf{Spatial logic.} The light source positions, blocking object
                      stop positions, and sun state that solve the puzzle are hidden;
                      players must deduce the correct configuration through observation
                      and reasoning rather than trial and error.
                \item \textbf{Shadow interdependence.} Moving a light source or blocking
                      object continuously reshapes the entire shadow pattern. A position
                      that creates one bridge may simultaneously remove another the player
                      needs. Players must reason about the full shadow layout, not just
                      local effects.
                \item \textbf{Trajectory constraints.} Blocking objects travel only along
                      their fixed paths and stop only at designer-specified positions.
                      The puzzle solution must be achievable within those constraints;
                      players cannot place objects arbitrarily.
            \end{itemize}

        \subsubsection{Appeal}
        Umbra appeals to fans of pure puzzle games---players who want a single elegant
        mechanic taken to its logical depth rather than narrative or action content.
        The two foundational rules combine to produce a surprisingly large puzzle space
        without requiring the player to learn many systems, making the game approachable
        to newcomers while offering genuine depth for dedicated puzzle solvers.

        Every level delivers a clean ``aha'' moment when the correct light arrangement
        snaps into place and the path becomes obvious. Solutions feel \emph{discovered},
        not stumbled upon---the satisfaction comes from reading the environment and
        inferring the hidden configuration, not from reaction speed or resource
        management.

        The game is accessible to players aged 12~and~up and structured for
        pick-up-and-play sessions of 5--15 minutes per level. It targets fans of
        \textit{The Witness}, \textit{Portal}, and \textit{Stephen's Sausage Roll}---games
        that trust the player to discover rules through play and reward logical deduction
        over reflex. The absence of combat, timers, and randomness makes Umbra equally
        welcoming to casual puzzle enthusiasts and genre veterans.

        \subsubsection{Look and Feel}
        Umbra uses a \textbf{minimalist silhouette aesthetic}. The scene is rendered with
        stark contrast between two dominant states: warm \textbf{amber} wash for areas in
        direct light and cool \textbf{blue-grey} for areas in shadow. Because shadows are
        continuous and computed in real time, the boundary between lit and shadowed
        regions is a smooth, live edge that shifts as the player drags a light source or
        adjusts a blocking object---making the cause-and-effect relationship between
        manipulation and shadow shape immediately visible.

        The visual language is deliberately simple so that the puzzle logic---not the
        art---demands the player's attention. Blocking objects are bold geometric
        silhouettes (squares, rectangles, cylinders) with no decorative surface detail.
        Light sources glow with a soft halo to distinguish them from passive objects.
        Draggable light sources show a subtle highlight ring while the mouse hovers over
        them; blocking objects show directional arrows along their trajectory to indicate
        available move directions. The player character is a minimal dark figure,
        readable against both lit and shadowed backgrounds.

        Ambient audio shifts in warmth with the sun state---brighter, airier tones for
        morning light; cooler, lower drones for evening shadow. Interaction feedback is
        delivered through crisp sound cues (light source drag start/end, blocking object
        slide, sun toggle) and the continuously animating shadow boundary, which gives
        the player instant visual confirmation of every manipulation.

    \subsection{Worlds, Characters and Story}

        \subsubsection{Back Story}
        % TODO: Provide any narrative framing or world-building context.

        \subsubsection{Spaces / Worlds}
        % TODO: Describe the level environments, tile types, light/shadow regions,
        %       and how they change based on sun state.

        \subsubsection{Characters}
        % TODO: Describe the player character and any non-player entities.

        \subsubsection{Levels of Difficulty}
        % TODO: Describe how puzzle complexity scales across levels.

    \subsection{Interaction Models}

        \subsubsection{User Interface --- Navigation and Movement}
        % TODO: Describe controls (mouse/keyboard, gamepad), camera perspective,
        %       and how the player moves through the world.

        \subsubsection{Game Play Sequence and Levels}
        % TODO: Describe the flow from level start to completion, including the
        %       morning/evening toggle interaction and win/lose conditions.

        \subsubsection{User/Environment --- Obstacles and Props}
        % TODO: Describe interactable objects (lanterns, spotlights, crates,
        %       pillars) and environmental hazards.

        \subsubsection{User/Character}
        % TODO: Describe how the player interacts with the player character
        %       (picking up objects, toggling sun state, etc.).

        \subsubsection{Character/Character}
        % TODO: Describe any character-to-character interactions (if applicable).

    \subsection{Performance and Scoring}

        \subsubsection{State Variables}
        % TODO: List game state variables tracked at runtime (sun state,
        %       object positions, level completion flag, move count, etc.).

        \subsubsection{Feedback}
        % TODO: Describe visual/audio feedback given to the player (shadow
        %       updates, success/failure cues, hints, etc.).

        \subsubsection{Performance and Progress Metrics}
        % TODO: Describe how player performance is measured and displayed
        %       (move count, time, star rating, level select, etc.).

% ─── 3. Implementation ───────────────────────────────────────────────────────
\section{Implementation}

    \subsection{Design Assumptions}

        \subsubsection{Hardware}
        % TODO: List minimum and recommended hardware requirements (CPU, RAM, GPU).

        \subsubsection{Software}
        % TODO: List software dependencies: Unreal Engine 5 version, OS (Windows),
        %       development tools, version control, etc.

        \subsubsection{Algorithms and Techniques}
        % TODO: Describe key technical approaches: Unreal Engine 5 real-time dynamic
        %       shadow system, continuous shadow geometry as physical collision, mouse
        %       drag interaction with boundary constraints for light sources, fixed
        %       trajectory/snap logic for blocking objects, sun-angle toggle, etc.

    \subsection{Storyboards}
    % TODO: Insert concept sketches or storyboard panels showing key game states
    %       (level start, object placement, sun toggle, level completion).

    \subsection{Design Logic}

        \subsubsection{Finite State Machines --- Events / States / Effects}
        % TODO: Provide FSM diagrams or tables for the game's core states:
        %       Idle, Object Selected, Sun Toggle, Level Complete, etc.

        \subsubsection{User Solution / Actions}
        % TODO: Walk through an example puzzle solution step-by-step, mapping
        %       player actions to state changes.

    \subsection{Software Versions}

        \subsubsection{Alpha Version}
        % TODO: List features included in the Alpha (vertical slice):
        %       one playable level, core shadow mechanic, sun toggle, basic UI.

        \subsubsection{Beta Version}
        % TODO: List additional features for Beta: multiple levels, polished
        %       art/audio, hint system, level select screen.

        \subsubsection{Tutorials}
        % TODO: Describe the in-game tutorial or introductory level that teaches
        %       the shadow mechanic and sun toggle to new players.

% ─── 4. Work Plan ────────────────────────────────────────────────────────────
\section{Work Plan}

    \subsection{Tasks}
    % TODO: Provide a numbered task/subtask breakdown with assigned team members
    %       and estimated durations. Example structure:
    %
    %  1.  Shadow/Light System
    %      1.1 Real-time dynamic shadow rendering (UE5)
    %      1.2 Shadow geometry as physical collision volumes
    %      1.3 Sun-toggle logic
    %  2.  Level Design
    %      2.1 Puzzle layout design (3 levels for alpha)
    %      2.2 Object placement editor
    %  3.  Player Interaction
    %      3.1 Object pickup/drop
    %      3.2 Sun toggle input
    %  4.  UI & HUD
    %      4.1 Level select screen
    %      4.2 Move counter / timer
    %  5.  Art & Audio
    %      5.1 Tile art (lit / shadow)
    %      5.2 Object models
    %      5.3 Sound effects & music
    %  6.  Testing & Polish

    \subsection{Milestones}

        \subsubsection{Minor Milestones}
        % TODO: List internal checkpoints with target dates.

        \subsubsection{Major Milestones}
        % TODO: List the course-defined major milestones:
        %  - Mon. Mar. 2:  Game Design Doc
        %  - Wed. Mar. 18: Alpha Version / Playtest
        %  - Wed. Mar. 25: Beta Version / Playtest
        %  - Mon. Mar. 30: Final Version / Demo
        %  - Wed. Apr. 1:  Game Trailer and Gameplay Videos

    \subsection{Development Schedule}
    % TODO: Insert a Gantt chart mapping Work Plan tasks to calendar weeks.
    %       The pgfgantt package is loaded above for this purpose.

\end{document}
